% TOP_PROBLEM: {"id": 1517, "title": "Hardy-Littlewood Conjecture A (Prime k-tuples)", "category_id": 1, "category": {"id": 1, "name": "number_theory", "display_name": "Number Theory", "description": "Properties of integers, prime numbers, Diophantine equations.", "slug": "number-theory", "order_index": 1, "created_at": "2026-07-31T15:26:25.670Z"}, "status": "open"}
% FIRST_POSTED: null
% ENTRY_KIND: statement_only
% Imported from: batch08.tex; UnsolvedMath ID 1517
\documentclass[11pt]{article}
\usepackage[margin=25mm]{geometry}
\usepackage{fontspec}
\setmainfont{Latin Modern Roman}
\usepackage{amsmath,amssymb,mathtools}
\usepackage{unicode-math}
\setmathfont{Latin Modern Math}
\usepackage{xurl,hyperref}
\hypersetup{colorlinks=true,urlcolor=blue}
\setcounter{secnumdepth}{0}
\setlength{\parindent}{0pt}
\setlength{\parskip}{5pt}
\setlength{\emergencystretch}{3em}
\newcommand{\sref}[2]{\hyperlink{source:#1:#2}{[S#2]}}
\newcommand{\cataloguescope}[1]{\paragraph{Recorded target.}#1}
\begin{document}
% UnsolvedMath ID 1517; source catalogue code HL-A
\setcounter{section}{1516}
\section{Hardy-Littlewood Conjecture A (Prime k-tuples)}\label{1517}
{\small\sffamily UnsolvedMath ID 1517 · Number Theory}
\subsection{Definitions and mathematical statement}

\paragraph{Shared notation.}$\mathbb N=\{1,2,\ldots\}$, and $\mathbb Z,\mathbb Q,\mathbb R,\mathbb C$ denote the integers, rationals, reals and complex numbers. Any other conventions are specified locally.


Let \(k\ge1\) and let \(a_1,\ldots,a_k\in\mathbb Z\) be fixed
offsets; repeated offsets may be deleted without changing the prime
conditions. The set of offsets is admissible if, for every prime \(p\),
there is some residue \(m\pmod p\) for which
\[
                  \gcd(m+a_i,p)=1\qquad(1\le i\le k).
\]
Equivalently, the residues \(-a_i\pmod p\) do not cover all residues
modulo \(p\). A prime is a positive integer larger than one with
exactly two positive divisors.
\paragraph{Statement recorded in the supplied source.}
For every fixed admissible collection of integer offsets
\(a_1,\ldots,a_k\), are there infinitely many \(n\in\mathbb N\) such that
\[
                         n+a_1,\ldots,n+a_k
\]
are all prime? Equivalently, does such an \(n\) exist above every
prescribed bound? \sref{1517}{2}
\subsection{Short English statement}
Does every fixed prime pattern without a modular obstruction occur infinitely often among translates of the integers?
\subsection{Sources}
\begin{description}
\item[\hypertarget{source:1517:1}{S1}] ULAM.AI and the UnsolvedMath Contributors, \emph{UnsolvedMath: A Curated Collection of Open Mathematics Problems}, record 1517 (HL-A). CC BY 4.0. \url{https://huggingface.co/datasets/ulamai/UnsolvedMath}.
\item[\hypertarget{source:1517:2}{S2}] James Maynard, \emph{Small gaps between primes}, Annals of Mathematics 181 (2015), 383--413; arXiv:1311.4600, Section 1, admissibility and the qualitative Prime $k$-tuples Conjecture. \url{https://arxiv.org/abs/1311.4600}.
\end{description}
\label{end:1517}
\end{document}
